\begin{explanationbox}
\textbf{为什么函数要求一个 $x$ 只能对应一个 $y$？}

\textbf{一、输入确定后，输出必须唯一。}
函数描述的是确定的对应规则：对定义域内每一个 $x$，都必须有且只有一个输出 $f(x)$。这样写出 $f(1)$，才知道它指哪一个数。这并不要求不同输入产生不同输出。
\begin{center}\begin{tikzpicture}[>=stealth,font=\small]
\draw[gray!40,rounded corners] (0,0) rectangle (14,7.5);
\node at (7,7) {$y=x^2$：每个输入只对应一个输出，允许不同输入同输出};
\begin{scope}[shift={(4.5,1.5)},x=1.65cm,y=1.1cm]
\draw[->,gray] (-1.8,0)--(1.9,0) node[right]{$x$};
\draw[->,gray] (0,-.3)--(0,3.7) node[above]{$y$};
\draw[blue,very thick,domain=-1.75:1.75,samples=100] plot(\x,{\x*\x});
\draw[orange!80!black,dashed] (-1,-.2)--(-1,3.4);
\draw[orange!80!black,dashed] (1,-.2)--(1,3.4);
\draw[gray,dotted] (-1,1)--(1,1);
\fill (-1,1) circle(2pt) node[left=8pt]{$(-1,1)$};
\fill (1,1) circle(2pt) node[right=8pt]{$(1,1)$};
\node[below] at (-1,0){$-1$};\node[below] at (1,0){$1$};
\end{scope}
\node[align=left,anchor=west] at (8.6,4.1) {$f(-1)=f(1)=1$\\[6pt]两条不同的竖线\\[4pt]各与图像交于一点};
\node at (7,.5) {关于 $y$ 轴对称：$(x,y)\longmapsto(-x,y)$，改变的是输入};
\end{tikzpicture}
\end{center}
图中 $-1$ 和 $1$ 都对应 $1$，完全符合函数定义。用竖直线检查图像：横坐标属于定义域时，竖线必须恰好交于一点；定义域外可以没有交点。

\newpage
\textbf{二、曲线方程不一定能表示 $y$ 关于 $x$ 的函数。}
\begin{center}\begin{tikzpicture}[>=stealth,font=\small]
\draw[gray!40,rounded corners] (0,0) rectangle (14,8);
\node at (7,7.5) {$y^2=4x$：同一条竖线交于两点，整条曲线不能写成 $y=f(x)$};
\begin{scope}[shift={(2.2,3.9)},x=1.4cm,y=1.05cm]
\draw[->,gray] (-.4,0)--(3.2,0) node[right]{$x$};
\draw[->,gray] (0,-2.8)--(0,2.9) node[above]{$y$};
\draw[blue,very thick,domain=0:2.6,samples=100] plot ({\x*\x/4},\x);
\draw[teal!80!black,very thick,domain=-2.6:0,samples=100] plot ({\x*\x/4},\x);
\draw[orange!80!black,dashed] (1,-2.7)--(1,2.7);
\fill (1,2) circle(2pt) node[right=5pt]{$(1,2)$};
\fill (1,-2) circle(2pt) node[right=5pt]{$(1,-2)$};
\node[below right] at (1,0){$x=1$};
\end{scope}
\node[anchor=west,align=left] at (7.6,4.3) {上支：$y=2\sqrt{x}$\\[8pt]下支：$y=-2\sqrt{x}$\\[8pt]各自的定义域都是 $[0,+\infty)$};
\node at (7,.5) {分开看，每一支都是 $y$ 关于 $x$ 的函数；合在一起则不是};
\end{tikzpicture}
\end{center}
一般地，$y^2=2px$（$p>0$）在每个 $x>0$ 处给出两个输出 $y=\pm\sqrt{2px}$；虽然 $x=0$ 处只有一个输出，仍不能使整条曲线成为 $y=f(x)$。限制为上支或下支后，各自才是函数。若改把 $y$ 当输入，则 $x=y^2/(2p)$ 又是 $x$ 关于 $y$ 的函数。

\newpage
\textbf{三、关于 $y$ 轴对称，为什么不违反函数定义？}
反射 $(x,y)\mapsto(-x,y)$ 配对的是两个输入，允许输出相同，因此偶函数可以关于 $y$ 轴对称。反射 $(x,y)\mapsto(x,-y)$ 却保持输入不变：只要 $y\ne0$，就会给同一个 $x$ 两个不同输出。
\begin{center}\begin{tikzpicture}[>=stealth,font=\small]
\draw[gray!40,rounded corners] (0,0) rectangle (14,6.8);
\node at (7,6.3) {关于 $x$ 轴对称：非零输出会冲突，零输出反射后仍是自己};
\begin{scope}[shift={(3.1,3.2)},x=1.3cm,y=1cm]
\draw[->,gray] (-1.5,0)--(2.2,0) node[right]{$x$};
\draw[->,gray] (0,-1.6)--(0,1.8) node[above]{$y$};
\draw[orange!80!black,dashed] (1,-1.5)--(1,1.5);
\fill[blue] (1,1) circle(2pt) node[above left]{$(1,1)$};
\fill[red!70!black] (1,-1) circle(2pt) node[below left]{$(1,-1)$};
\node[align=center] at (.3,-2.3) {若同时保留这两个点\\同一输入 $1$ 就有两个输出};
\end{scope}
\begin{scope}[shift={(10.3,3.2)},x=1.3cm,y=1cm]
\draw[->,gray] (-1.7,0)--(2,0) node[right]{$x$};
\draw[->,gray] (0,-1.6)--(0,1.8) node[above]{$y$};
\draw[blue,very thick] (-1.6,0)--(1.7,0);
\draw[orange!80!black,dashed] (1,-1.4)--(1,1.5);
\fill[orange!80!black] (1,0) circle(2pt) node[above left]{$(1,0)$};
\node at (-.95,1.3){$f(x)\equiv0$};
\node[align=center] at (0,-2.3) {镜像点与原来的点重合\\同一输入仍只有输出 $0$};
\end{scope}
\end{tikzpicture}
\end{center}
所以，非空定义域上的实值函数图像若关于 $x$ 轴对称，必有 $f(x)=-f(x)$，即定义域内处处 $f(x)=0$；反过来，处处为零的函数图像确实关于 $x$ 轴对称。注意：它要同时为奇函数和偶函数，还需定义域关于原点对称。

\medskip
\textbf{回到奇偶性：}先保证每个输入有唯一输出，再检查 $x$ 与 $-x$ 是否都在定义域内，最后比较 $f(-x)$ 与 $f(x)$。
\end{explanationbox}
\newpage
\begin{explanationbox}
	\textbf{\textcolor{CExplanation}{一句话直觉}}

	函数的奇偶性本质上是对称性：奇函数关于原点对称，偶函数关于$y$轴对称。要谈对称，定义域必须关于原点对称才有意义。

	\medskip
	\textbf{\textcolor{CExplanation}{核心拆解}}
	\begin{description}[style=nextline, leftmargin=2.8em, labelsep=0.5em, itemsep=0.45em]
		\item[\textbf{要点一（定义域对称）：}] 若定义域不关于原点对称，则函数一定不具有奇偶性。这是前提条件。
		\item[\textbf{要点二（关系式检验）：}] 在定义域对称的前提下，计算$f(-x)$并与$f(x)$比较：相等则为偶函数，相反则为奇函数。
	\end{description}

	\medskip
	\textbf{\textcolor{CExplanation}{几何本质（对称变换）}}

	奇函数：图像绕原点旋转$180^\circ$后与自身重合。偶函数：图像关于$y$轴作镜面对称后与自身重合。

	\medskip
	\textbf{\textcolor{CExplanation}{代数意义（简化运算）}}

	奇偶性可减少计算量：例如偶函数满足$f(x)=f(|x|)$，奇函数满足$f(-x)=-f(x)$，因此只需研究$x\ge0$部分。

	\medskip
	\textbf{\textcolor{CExplanation}{考点价值（判断与运用）}}
	\begin{itemize}[leftmargin=*, itemsep=0.5em, topsep=0.4em]
		\item \textbf{考法一（奇偶性判定）：} 给定函数表达式，要求判断奇偶性。必须先检查定义域对称，再计算$f(-x)$。
		\item \textbf{考法二（利用奇偶求参）：} 已知函数为奇函数或偶函数，求参数值。通常利用特殊值法（如$f(0)=0$仅对奇函数有效）或恒等式法。
	\end{itemize}

	\medskip
	\textbf{\textcolor{CExplanation}{顿悟点}}

	先看定义域！很多学生做题第一步就计算$f(-x)$，忽略定义域对称性，导致判断错误。

	\medskip
	\textbf{\textcolor{CExplanation}{使用场景}}
	\begin{itemize}[leftmargin=*, itemsep=0.5em, topsep=0.4em]
		\item \textbf{场景一（直接判定）：} 遇到具体函数，按步骤：定义域$\to$ $f(-x)$ $\to$结论。
		\item \textbf{场景二（奇偶应用）：} 已知函数奇偶性，利用其性质简化对称区间上的运算或方程求解。
	\end{itemize}
\end{explanationbox}
